\documentclass{situationist-preprint}

\title{On the Meaning of Life}
\msid{EXI-01-a}
\msfamily{EXI-01}
\msdiscipline{Meaning, Death and the Absurd}
\mstype{Principal manuscript}
\mslean{lean/Philosophy/EXI/EXI\_01.lean}

\begin{document}
\maketitle
\vspace{-1em}

\noindent\textbf{Definition 1.} A life $L$ is \emph{meaningful to degree $d$} if $d$ is the extent to which $L$ consists in active engagement with projects of objective worth \citep{Wolf2010}.

\medskip\noindent\textbf{Assumption 2.} There exist projects of objective worth: projects whose value does not depend on the attitudes of those engaged in them \citep{Metz2013}.

\medskip\noindent\textbf{Assumption 3.} For every agent capable of intentional action, there is at least one project of objective worth in which she is able to engage.

\medskip\noindent\textbf{Theorem 4.} For every such agent, a life of positive meaning is available.

\medskip\noindent\textit{Proof.} By Assumption 3, there is a project $p$ of objective worth in which the agent can engage. A life that includes engagement with $p$ consists, to some positive extent, in active engagement with a project of objective worth. By Definition 1, it is meaningful to a positive degree. \hfill$\square$

\medskip\noindent\textbf{Remark 5.} The theorem answers the question of what the meaning of life is, and of whether meaning is available to everyone. It does not answer the question of whether any particular life is meaningful, which depends on what is done with it. Nor does it answer the sceptical challenge that, \emph{sub specie aeternitatis}, nothing matters \citep{Nagel1971}: Assumption 2 denies the premise of that challenge rather than refuting it. A defence of Assumption 2, and a discussion of the absurd, are given in the commentary (EXI-01-b).

\begin{thebibliography}{9}\small
\bibitem[Metz(2013)]{Metz2013} Metz, T. (2013). \emph{Meaning in Life: An Analytic Study}. Oxford University Press.
\bibitem[Nagel(1971)]{Nagel1971} Nagel, T. (1971). The absurd. \emph{The Journal of Philosophy}, 68(20), 716--727.
\bibitem[Wolf(2010)]{Wolf2010} Wolf, S. (2010). \emph{Meaning in Life and Why It Matters}. Princeton University Press.
\end{thebibliography}

\end{document}
